\chapter{Summary and Prospects}
\label{ch:summary}

\section{Summary}
\label{sec:summary}

This thesis set out to migrate a Nav2-based navigation stack for an
Ackermann-steered vehicle, the Quanser QCar, from a validated Gazebo
simulation onto the physical robot, and to identify and address the
sim-to-real gap encountered along the way
(Section~\ref{sec:problem-statement}). Each of the four sub-questions
posed in the introduction is answered directly by one methodology
chapter:

\begin{itemize}
  \item \textbf{Simulation-stack structure} (Chapter~\ref{ch:sim-stack}):
    the stack was migrated twice -- first from a hand-rolled bicycle-model
    inversion to \texttt{ros2\_control}'s Ackermann controller
    (deliberately adopting real dead-reckoning odometry rather than
    simulation ground truth, so the eventual hardware migration would
    not have to absorb that discrepancy on top of everything else),
    then to a direct Gazebo drive plugin -- which, as
    Section~\ref{sec:odom-regression} documents, quietly reverted
    \texttt{/odom} back to ground truth despite the topic name and
    intent staying the same, a discrepancy this thesis surfaces rather
    than leaves uncorrected.
  \item \textbf{Bridging incompatible ROS~2 distributions}
    (Chapter~\ref{ch:hardware-bridge}): native ROS~2 communication
    between the vehicle's onboard ROS~2 Dashing and the development
    host's ROS~2 Humble was confirmed, through direct testing rather
    than assumption, to be impossible at the DDS protocol level. A
    TCP/JSON bridge was designed and built that presents the physical
    robot to the rest of the navigation stack as an ordinary ROS~2
    robot, with LiDAR angle-convention, odometry, steering-trim, and
    inter-machine clock-synchronisation corrections layered on top to
    bring the bridged data to a fidelity usable by AMCL and Nav2.
  \item \textbf{Actuator identification and calibration}
    (Chapters~\ref{ch:hardware-bridge}--\ref{ch:calibration}): an
    uncalibrated throttle mapping (2.6--2.7$\times$ overshoot) was
    corrected to within 5\% accuracy across
    \SIrange{0.3}{0.5}{\meter\per\second} via an empirical duty sweep;
    a large ($\approx$\SI{2}{\second}) apparent coasting delay was
    ultimately traced not to physical drivetrain dynamics but to a
    ring-buffer artefact in the vendor hardware abstraction layer's
    task-based I/O mode, reducing stiction delay and stop latency by
    more than an order of magnitude once corrected; and, independently,
    a real $\approx$39\% steering-gain error -- separate from, and not
    fixed by, the vehicle's zero-point trim -- was characterised
    through a seven-point sweep and corrected to $\approx$5.6\%
    residual error (Section~\ref{sec:steering-gain}).
  \item \textbf{Closed-loop validation} (Chapter~\ref{ch:nav-integration}):
    the bridged, calibrated vehicle was shown to complete real Nav2
    goals requiring both straight driving and turning, using a custom
    Ackermann-appropriate MPPI critic and a phased integration
    approach that isolated each subsystem before combining it with the
    next. A separate execution-side freeze at reversal cusps during
    K-turn manoeuvres in tight spaces was diagnosed to MPPI's own
    cusp-arrival gating and resolved with a custom path-smoothing
    plugin, cutting a representative manoeuvre's mean duration from
    \SI{110.2}{\second} (with recoveries) to \SI{43.9}{\second} (with
    none) in simulation, and validated directly on the physical vehicle
    across two independent goals with zero recoveries in either
    (Section~\ref{sec:cusp-freeze}).
\end{itemize}

Chapter~\ref{ch:results} quantified the resulting gap: actuator
response now tracks commanded speed within 5\% (down from a
2.6--2.7$\times$ error), and stop-latency/coast behaviour was reduced
from roughly 2.4\,s and up to 3$\times$ the commanded stop distance to
approximately 0.5\,s and $\approx$\SI{0.1}{\meter}, a residual fully
explained by ordinary linear velocity decay rather than any remaining
unexplained mechanism.

\subsection{Critical Reflection}
\label{sec:critical-reflection}

What worked well in this project was the discipline of measuring
directly on hardware rather than trusting a plausible-sounding
explanation: four of this thesis's most significant findings (the
protocol-version incompatibility in Section~\ref{sec:rmw-incompatibility},
the sim-vs-real overshoot ratio in Section~\ref{sec:sim-vs-real}, the
buffering defect in Section~\ref{sec:readmode-bug}, and the
reversal-cusp freeze in Section~\ref{sec:cusp-freeze}) were each, at
some point, plausibly attributable to a simpler cause -- localisation
error, a request-delivery bug, physical coasting, a tuning-only
gain problem -- that direct measurement subsequently ruled out. The
cusp-freeze investigation additionally illustrates the same discipline
applied to solutions, not only diagnoses: three plausible fix attempts
and two later follow-on refinements were each tested directly against
hardware or a controlled reproduction, found wanting by measurement
rather than intuition, and reverted with the evidence that ruled them
out recorded rather than discarded. What could have gone better is
equally clear from Section~\ref{sec:active-braking}: a reactive
control design was deployed to physical hardware before it had been
adequately traced through or simulated, and it produced a genuine
safety incident. That the incident caused no lasting harm does not
make the shortcut it resulted from acceptable practice, and this
thesis records it, unedited, as a documented counter-example to its
own otherwise measurement-first methodology.

A second, more structural limitation is that not every quantitative
result in this thesis is a genuinely final, post-fix measurement.
Section~\ref{sec:display-lag-interpretation} reports navigation
timing figures measured while the buffering defect of
Section~\ref{sec:readmode-bug} was still active, with only a
qualitative (not re-instrumented) confirmation available after the
fix, and Section~\ref{sec:results-sim-vs-real}'s sim-versus-real
comparison likewise predates both major fixes in
Chapter~\ref{ch:calibration}. These are reported as open items rather
than papered over, and the exact next steps to close them --
re-running \texttt{monitor\_goal\_timing.py} and the six-point
sim-versus-real sweep under the final, fixed configuration -- are
specific and immediately actionable rather than vague.

\section{Prospects}
\label{sec:prospects}

Beyond the two re-measurements identified above, three directions
follow naturally from this thesis's findings:

\begin{itemize}
  \item \textbf{Closed-loop low-speed control.} This thesis's
    calibration (Section~\ref{sec:final-calibration}) explicitly
    identified the sub-\SI{0.25}{\meter\per\second} regime as a hard
    limit of open-loop duty control, not a remaining tuning gap.
    Reliable low-speed creep -- useful for close-quarters manoeuvring
    or fine final-approach positioning -- would require closed-loop
    velocity feedback from the encoder rather than the fixed duty
    mapping used throughout this thesis.
  \item \textbf{Deceleration-profile tuning for coast reduction.} The
    stop-latency residual characterised in Section~\ref{sec:results-stop-distance}
    is now small, but not zero, and scales with cruise speed. Given
    that active braking is explicitly ruled out on safety grounds
    (Section~\ref{sec:active-braking}), a lower-risk avenue --
    identified but not investigated within this thesis -- is tuning
    the MPPI \texttt{GoalCritic}/deceleration profile in
    \texttt{nav2\_params.yaml} so the controller commands a lower
    cruise speed during final approach to a goal, proportionally
    shrinking (not eliminating) the coast without introducing any new
    reactive control loop.
  \item \textbf{Generalising the diagnostic methodology.} This
    thesis's most transferable contribution may be procedural rather
    than numeric: the recurring pattern of an internally
    self-consistent software abstraction masking a stale or defective
    input (Section~\ref{sec:readmode-bug}), and the specific
    discipline of verifying a component's own upstream input rather
    than only its internal consistency, is a general principle worth
    testing on other buffered or task-based hardware interfaces beyond
    the one vendor library examined here.
  \item \textbf{Closing the simulated-versus-real replan-complexity gap.}
    Section~\ref{sec:cusp-hw-validation} found that real-hardware K-turn
    manoeuvres occasionally produce substantially higher per-replan
    cusp counts (8--10) than any simulation trial reproduced (maximum
    $\approx$3), without preventing either tested goal from succeeding.
    The leading candidate explanations -- \texttt{SmacPlannerHybrid}'s
    sensitivity to real, more awkward goal/environment geometry, or a
    residual, session-specific clock-offset effect similar to the one
    implicated in Section~\ref{sec:cusp-followon-reverted}'s reverted
    \texttt{IsPathValidDebounced} attempt -- were identified but not
    isolated, and remain open. Relatedly, a deviation-aware critic
    weighting scheme -- easing a sampled trajectory's pull back toward
    the exact planned path once it has deviated significantly, in
    favour of a goal-direction target, rather than insisting on precise
    reconvergence -- was proposed but not implemented, as a
    controller-side complement to Section~\ref{sec:cusp-freeze}'s
    path-preprocessing fix.
  \item \textbf{Hardware re-validation of the curvature-continuing
    smoother refinement.} Section~\ref{sec:cusp-continued-iteration}'s
    two most recent refinements -- the \texttt{min\_lead\_distance} fix
    to the cusp re-extension bug, and the curvature-continuing
    extension itself -- were, at the time of writing, verified only in
    simulation (a single smoke test for the curvature refinement, two
    trials without directly observing the skip branch fire for
    \texttt{min\_lead\_distance}). Both are strict generalisations of
    the already hardware-validated straight-line mechanism and are
    expected to transfer, but this thesis deliberately does not claim
    hardware confirmation it does not have; repeating
    Section~\ref{sec:cusp-hw-validation}'s two-goal real-vehicle test
    against this refined configuration is the immediate next step.
\end{itemize}

The QCar's navigation stack, as delivered by this thesis, reaches
real closed-loop navigation goals on physical hardware with a
sim-to-real gap that is now measured, explained, and in most respects
substantially reduced -- the original goal stated in
Section~\ref{sec:problem-statement}. What remains open is narrower
and better characterised than what this thesis began with: not
``why doesn't this work,'' but two specific, already-scoped
measurements and two specific, already-identified directions for
further improvement.
